\documentclass[11pt,a4paper]{article}
\usepackage[T1]{fontenc}
\usepackage{lmodern}
\usepackage[margin=25mm]{geometry}
\usepackage{amsmath,amssymb,amsthm,mathtools}
\usepackage{booktabs,longtable,array}
\usepackage{xcolor,listings}
\usepackage[hidelinks]{hyperref}
\allowdisplaybreaks
\setlength{\emergencystretch}{2em}
\newtheorem{theorem}{Theorem}[section]
\newtheorem{lemma}[theorem]{Lemma}
\newtheorem{proposition}[theorem]{Proposition}
\newtheorem{remark}[theorem]{Remark}
\newcommand{\dd}{\mathrm d}
\newcommand{\tr}{\mathsf T}
\newcommand{\NF}{\operatorname{NF}}
\newcommand{\Res}{\operatorname{Res}}
\newcommand{\diag}{\operatorname{diag}}
\newcommand{\e}{\mathbf e}
\newcommand{\vv}{\mathbf v}
\newcommand{\ww}{\mathbf w}
\newcommand{\RR}{\mathbb R}
\newcommand{\CC}{\mathbb C}
\newcommand{\QQ}{\mathbb Q}
\lstset{basicstyle=\ttfamily\footnotesize,breaklines=true,
 columns=fullflexible,keepspaces=true,showstringspaces=false,
 frame=single,rulecolor=\color{black!25},tabsize=4}

\title{A rational contour proof of a Ramanujan-type\texorpdfstring{ $1/\pi^2$}{ 1/pi squared} identity}
\author{Hodge project: scalar Row12 manuscript}
\date{7 October 2026}

\begin{document}
\maketitle
\begin{abstract}
We evaluate a specific convergent weighted factorial series. The proof first
converts the entire series into an angular and radial integral. An explicit
rational coefficient one-form and a six-component rational correction give
an antiderivative of the loaded radial integrand. Both endpoint limits and
the removable interior singularity are treated explicitly. The nonzero
endpoint is evaluated using an independently normalized Gauss Wronskian and
a signed algebraic transformation. All finite rational inputs and an exact
verification procedure are included in this single source file.
\end{abstract}

\noindent\textbf{Version and verification scope.}
This manuscript concerns the literal scalar equality below. Its analytic
derivation has evidence tags [ARG], and the retained exact rational checks
have tag [ONE]. A fresh document-level review found no substantive defect
in the scalar proof; its exact scope and source versions are recorded in
\texttt{DOCUMENT\_REVIEW.md}. Human mathematical acceptance is Pending.
The literal scalar identity also has a completed Lean formalization as of
7 October 2026, including unconditional convergence and sum statements.
\tableofcontents

\section{The statement and the proof structure}
\begin{theorem}\label{thm:main}
The series on the left converges absolutely, and
\begin{equation}\label{eq:main}
 R:=\sum_{n=0}^{\infty}\frac{(6n)!}{(n!)^6}
       \frac{532n^2+126n+9}{10^{6n}}
   =\frac{375}{4\pi^2}.
\end{equation}
\end{theorem}

Set
\begin{equation}\label{eq:constants}
 \rho=\frac{729}{15625}=\frac{6^6}{10^6},\qquad
 u=\frac{27}{125}=\sqrt\rho,\qquad \tau=2\pi i,
\end{equation}
and introduce the two distinct weight polynomials
\begin{equation}\label{eq:weights}
 P(z)=9+126z+532z^2,\qquad
 Q(z)=(1+6z)P(z)=9+180z+1288z^2+3192z^3.
\end{equation}
The factor $1+6z$ in $Q$ is essential: it comes from the radial moment.
The argument constructs a function $\mathcal E$ satisfying
\[
 \mathcal E'(U)=-\frac{1}{3375}
 Q(\theta_\lambda)h\bigl(\rho(1-U^2)^3\bigr),
 \qquad \theta_\lambda=\lambda\frac{\dd}{\dd\lambda},
\]
where $Q(\theta_\lambda)h$ is evaluated at the displayed argument.
We prove, including every boundary contribution,
\[
 R=\int_0^1Q(\theta_\lambda)h\bigl(\rho(1-U^2)^3\bigr)\,\dd U,
 \qquad \mathcal E(0+)=\frac{1}{36\pi^2},\quad
 \mathcal E(1-)=0.
\]
The interior rational pole of the initially defined $\mathcal E$ is removable.
These facts imply~\eqref{eq:main} by the fundamental theorem of calculus.

\section{From the literal series to the loaded integral}
For $n\geq0$ define
\begin{equation}\label{eq:native-coeff}
 a_n=\frac{(1/6)_n(1/2)_n(5/6)_n}{(n!)^3}
     =\frac{(6n)!}{(3n)!(n!)^3\,1728^n},
\end{equation}
where $(a)_0=1$ and $(a)_n=a(a+1)\cdots(a+n-1)$.
The equality of the two formulas follows from their common value $1$ at
$n=0$ and their common successive ratio
\begin{equation}\label{eq:native-ratio}
 \frac{a_{n+1}}{a_n}
 =\frac{(n+1/6)(n+1/2)(n+5/6)}{(n+1)^3}.
\end{equation}
In particular $0<a_n\leq1$. Define, on the complex unit disk,
\begin{equation}\label{eq:series}
 F(q)=\sum_{n\geq0}a_nq^n,\qquad
 h(\lambda)=\sum_{n\geq0}a_n^2\lambda^n,\qquad
 \theta=q\frac{\dd}{\dd q}.
\end{equation}
The variable attached to each Euler operator will always be stated.
All finite Euler derivatives converge normally on compact subsets of their
unit disks, by the majorant $\sum n^j r^n$ for $r<1$.

\begin{lemma}[Radial moments]\label{lem:moments}
For every integer $m\geq0$,
\[
 I_m:=\int_0^1(1-U^2)^m\,\dd U
     =\frac{4^m(m!)^2}{(2m+1)!}.
\]
Consequently
\begin{equation}\label{eq:moment-bridge}
 (6n+1)a_n^2 I_{3n}=
       \frac{(6n)!}{6^{6n}(n!)^6}.
\end{equation}
\end{lemma}
\begin{proof}
The derivative of $U(1-U^2)^{m+1}$ has integral zero: the function
vanishes at both endpoints. Its integral gives
$(2m+3)I_{m+1}=(2m+2)I_m$. Starting at $I_0=1$ proves the stated formula.
Substitution of~\eqref{eq:native-coeff} with $m=3n$ proves
\eqref{eq:moment-bridge}; here $4^{3n}/1728^{2n}=6^{-6n}$.
\end{proof}

Absolute convergence of the original series follows independently from
$(6n)!/(n!)^6\leq6^{6n}$, since this multinomial coefficient is one term
in the expansion of $(1+\cdots+1)^{6n}$ with six summands. Thus its terms
are bounded by $P(n)\rho^n$, a summable sequence.
Put $x=1-U^2$ and $\lambda=\rho x^3$. For $0\leq U\leq1$,
the loaded summands are bounded by $Q(n)\rho^n$. This justifies
termwise radial integration and yields the exact identity
\begin{equation}\label{eq:radial}
 R=\int_0^1Q(\theta_\lambda)h\bigl(\rho(1-U^2)^3\bigr)\,\dd U.
\end{equation}
Indeed the $n$th integrated term is $Q(n)a_n^2\rho^nI_{3n}$,
which equals the $n$th original term by~\eqref{eq:moment-bridge}.
The term $n=0$ has value $9$ in both expressions.

\begin{lemma}[Angular coefficient extraction]\label{lem:hadamard}
For $0<\lambda<1$ and $\lambda<r<1$, every positively oriented circle
$|s|=r$ satisfies
\begin{equation}\label{eq:hadamard}
 \theta_\lambda^j h(\lambda)
 =\frac1\tau\oint_{|s|=r}
       F(s)(\theta_b^j F)(\lambda/s)\,\frac{\dd s}{s},
 \qquad j\geq0.
\end{equation}
Here $\theta_b=b\,\dd/\dd b$, and for $j=0$ the factor is $F(b)$.
\end{lemma}
\begin{proof}
Both native power series and their required Euler derivatives converge
absolutely and uniformly on the circle. Expanding the product, the
normalized integral of $s^{m-n}\dd s/s$ equals $1$ when $m=n$ and
$0$ otherwise. This leaves $\sum a_n^2n^j\lambda^n$.
For $j=0$, the factor $n^0$ denotes $1$, including at $n=0$.
The bounds already given also justify local parameter differentiation.
\end{proof}

\section{The explicit rational coefficient form}
The coefficient recurrence gives the native differential equation
\[
 \bigl[\theta^3-q(\theta+1/6)(\theta+1/2)(\theta+5/6)\bigr]F=0.
\]
Writing $Y(q)=(F(q),\theta F(q),\theta^2 F(q))^{\tr}$, this becomes
\begin{equation}\label{eq:native-system}
 \theta Y=N(q)Y,\qquad
 N(q)=\begin{pmatrix}
 0&1&0\\0&0&1\\
 \dfrac{5q}{72(1-q)}&\dfrac{23q}{36(1-q)}&\dfrac{3q}{2(1-q)}
 \end{pmatrix}.
\end{equation}
Set
\begin{equation}\label{eq:dual-state}
 T(q)=\begin{pmatrix}
 -5q/72&0&2q^2(1-q)\\
 -q/2&q(q-1)&0\\1-q&0&0
 \end{pmatrix},\qquad W(q)=T(q)^{\tr}Y(q).
\end{equation}
For use in the finite coefficient calculation, let
\begin{equation}\label{eq:sym-connection}
 \mathsf A(q)=\begin{pmatrix}0&1&0\\2c(q)&-a(q)&2\\0&c(q)&-2a(q)\end{pmatrix},
 \quad a(q)=\frac{1-3q/2}{q(1-q)},\quad
 c(q)=\frac{5}{144q(1-q)}.
\end{equation}
Entrywise rational differentiation gives
$T'+N^{\tr}T/q+T\mathsf A=0$. Hence
\begin{equation}\label{eq:dual-ode}
 W'(q)=-\mathsf A(q)^{\tr}W(q).
\end{equation}
This identity fixes the dual-state convention used below.

For $0<U<1$ set
\begin{equation}\label{eq:beta-parameters}
 \begin{aligned}
 x&=1-U^2,& \lambda&=\rho x^3,& b&=\lambda/s,\\
 m&=\frac{27x(80-99x+27x^2)}{1000},&
 \ell_0&=\frac{27x(9x-16)}{1000},& \ell&=U\ell_0,\\
 M&=25-9x,& d&=x^4M(99x-50)^2,& f&=s^2-2ms+\lambda.
 \end{aligned}
\end{equation}
The polynomial matrices $L_P(x),R_P(x)$ are explicitly specified in
Appendix~\ref{app:certificates}. Define the rational matrix coefficients
\begin{align}
 B_s(U,s)&=\frac{2U L_P(s-m)+2\ell M R_P}{df},\label{eq:beta-s}\\
 B_U(U,s)&=\frac{U L_P\,\partial_U f
 -2R_P\{\ell M\partial_Um+(s-m)(M\partial_U\ell+9U\ell)\}}{df}.
 \label{eq:beta-u}
\end{align}
The partial $U$ derivatives hold $s$ fixed. Both coefficients are used;
the second summand in~\eqref{eq:beta-s} is retained throughout.
Introduce the actual normalized scalar contour function
\begin{equation}\label{eq:beta-flux}
 \mathfrak b(U)=\frac1\tau\oint_{|s|=\sqrt\lambda}
       W(s)^{\tr}B_s(U,s)W(\lambda/s)\,\dd s.
\end{equation}

The polynomial identity
$m^2-\lambda=U^2\ell_0^2M$ gives the two roots
\begin{equation}\label{eq:roots}
 \sigma_\pm=m\pm U\ell_0\sqrt M,
 \qquad 0<\sigma_+<\sqrt\lambda<\sigma_-.
\end{equation}
Indeed $m>0$, $\ell_0<0$, $M>0$ and
$\sigma_+\sigma_-=\lambda$ on $0<x<1$. The circle therefore contains
exactly the inside root $\sigma_+$ and has no pole on it. The coefficients
also have a rational pole at
\begin{equation}\label{eq:escape}
 99x-50=0,\qquad U_e=\frac{7}{3\sqrt{11}}.
\end{equation}
For now~\eqref{eq:beta-flux} is used off this point.

\section{The finite identities and the scalar primitive}
This section makes explicit how finite rational identities imply an
identity of actual analytic functions.
Let $J=N(0)$, $\ww=(5/72,23/36,3/2)$, and let $E$ be the $3$ by $3$
matrix whose last row is $\ww$ and whose other rows vanish. Write
\[
 \mathcal M_s=
 \frac{N(s)\otimes I-I\otimes N(\lambda/s)}s.
\]
Flatten $3$ by $3$ matrices row by row; $\e_0,\e_1,\e_2$ denote standard
rows of length three. Define the six rows
\begin{equation}\label{eq:nf-basis}
 \Psi_j=\frac{\ww\otimes\e_j}{s-1},\qquad
 \Psi_{j+3}=\frac{\e_j\otimes\ww}{s-\lambda},\qquad 0\leq j\leq2.
\end{equation}
The corresponding column of contour integrals is
\begin{equation}\label{eq:v-definition}
 v_j(\lambda)=\frac1\tau\oint_{|s|=r}
  \Psi_j(s,\lambda)\,[Y(s)\otimes Y(\lambda/s)]\,\dd s,
 \qquad \lambda<r<1.
\end{equation}
These integrals are independent of the admissible radius, since their
integrands are analytic in the intervening annulus. They are holomorphic
in $\lambda$ locally on the punctured unit disk by a fixed-circle argument:
for complex $\lambda$ choose $|\lambda|<r<1$ in~\eqref{eq:v-definition}.

\begin{lemma}[Finite rational certificate]\label{lem:certificate}
With the explicit constant matrix $C$, rational matrix $B_6(\lambda)$
and rational row $\xi(x)$ in Appendix~\ref{app:certificates}, one has
\begin{align}
 C\vv(\lambda)&=(h,\theta_\lambda h,\ldots,\theta_\lambda^5h)^{\tr},
 \label{eq:cyclic-state}\\
 \theta_\lambda\vv&=B_6(\lambda)\vv.\label{eq:six-action}
\end{align}
Define the two covariant coefficient derivatives
\begin{align*}
 \mathcal D_U B_s&=\partial_U B_s+\frac{6U}{x}B_s\,b\mathsf A(b)^{\tr},\\
 \mathcal D_s B_U&=\partial_s B_U-\mathsf A(s)B_U
                      +\frac{b}{s}B_U\mathsf A(b)^{\tr},
\end{align*}
and $\mathsf S=T(s)(-\mathcal D_U B_s+\mathcal D_sB_U)T(b)^{\tr}$.
There is a rational row $\eta$ with poles only at $0,1,\lambda$ and
infinity in $s$, such that
\begin{equation}\label{eq:source-nf}
 \operatorname{flatten}(\mathsf S)
   =\sum_{j=0}^5G_j(x)\Psi_j+\partial_s\eta+\eta\mathcal M_s.
\end{equation}
Its six source coordinates satisfy the complete row identity
\begin{equation}\label{eq:outer-certificate}
 G-\frac{9c_0+180c_1+1288c_2+3192c_3}{3375}
 =\xi-2(1-x)\xi'-\frac{6(1-x)}x\xi B_6(\rho x^3),
\end{equation}
where $c_j$ is row $j$ of $C$ (indices start at zero).
All statements are rational identities on their regular charts.
\end{lemma}
\begin{proof}
Appendix~\ref{app:certificates} specifies every rational input and the
finite elimination giving the primitive $\eta$. It verifies the anchor,
the six parameter-action identities, and~\eqref{eq:outer-certificate}
using exact rational arithmetic. These identities, rather than a rank or
numerical residual tolerance, are the finite certificate.
For an exact row $\partial_s\zeta+\zeta\mathcal M_s$, its pairing with
$Y(s)\otimes Y(\lambda/s)$ is the ordinary derivative of the single-valued
function $\zeta[Y(s)\otimes Y(\lambda/s)]$. Its integral on the circle is
zero. The anchor therefore gives $h=c_0\vv$, and the parameter-action
identities give~\eqref{eq:six-action} and the subsequent five rows of
\eqref{eq:cyclic-state}. No regularity inside the circle is required for
the integral of this total derivative to vanish.
\end{proof}

For any $U_0\in(0,1)\setminus\{U_e\}$, root separation in~\eqref{eq:roots}
allows a fixed radius in a neighborhood of $U_0$ without crossing a pole
of $f$. On this fixed compact circle all integrands and their derivatives
are analytic in the parameter. Applying~\eqref{eq:dual-ode} and
integrating the $s$ derivative of the $B_U$ pairing gives
\begin{equation}\label{eq:beta-derivative}
 \mathfrak b'(U)=-G(x)\vv(\lambda).
\end{equation}
Here~\eqref{eq:source-nf} supplies the source integral, with its displayed
minus sign. The coefficient derivatives account for
$\partial_Ub=-6Ub/x$ and $\partial_sb=-b/s$.

By~\eqref{eq:six-action}, ordinary differentiation gives
\[
 \frac{\dd}{\dd U}\{U\xi(x)\vv(\lambda)\}
 =\left[\xi-2(1-x)\xi'-\frac{6(1-x)}x\xi B_6\right]\vv.
\]
Combining this with~\eqref{eq:outer-certificate} and
\eqref{eq:beta-derivative}, define
\begin{equation}\label{eq:primitive}
 \boxed{\mathcal E(U)=\mathfrak b(U)+U\xi(1-U^2)\vv\bigl(\rho(1-U^2)^3\bigr).}
\end{equation}
It satisfies the actual function identity
\begin{equation}\label{eq:primitive-derivative}
 \mathcal E'(U)=-\frac1{3375}Q(\theta_\lambda)h(\lambda),
 \qquad \lambda=\rho(1-U^2)^3.
\end{equation}

\section{The interior escape singularity is removable}
At $x_e=50/99$, direct substitution gives the two roots $\lambda_e$ and
$1$ of $f$, with $0<\lambda_e<1$. Choose a fixed circle of radius
$\lambda_e<r_e<1$. For complex $U$ sufficiently close to $U_e$, its
integrands remain analytic on the circle. The continued integral and the
rational correction in~\eqref{eq:primitive} are meromorphic in $U$;
their only possible local pole comes from powers of $99x-50$.
For real nearby $U$ this circle is deformable to the balanced circle without
crossing either root, so it represents the same function $\mathcal E$.

The right-hand side of~\eqref{eq:primitive-derivative} is holomorphic near
$U_e$. Consequently $\mathcal E$ has no negative Laurent coefficient:
the derivative of a nonzero term $a_{-k}(U-U_e)^{-k}$ would have a pole
of order $k+1$. Thus $\mathcal E$ extends holomorphically through $U_e$.
Both real one-sided limits agree. This conclusion concerns the sum of the
two terms in~\eqref{eq:primitive}; every rational escape term in each
coefficient remains included.

\section{The upper endpoint and the rational correction}
\subsection{The beta contour tends to zero at \texorpdfstring{$U=1$}{U=1}}
As $q\to0$ in a disk,
$F(q)=1+(5/72)q+O(q^2)$. Formula~\eqref{eq:dual-state} shows that all
three components of $W(q)$ are $O(q^2)$. In the first component the
potentially linear terms are $-5q/72+(5/72)q$ and cancel.
The bounds are uniform on a sufficiently small disk.

As $U\to1^-$, put $r=\sqrt\lambda=u x^{3/2}$. On $|s|=r$ both $s$
and $b=\lambda/s$ have modulus $r$, so the product of the two dual-state
bounds is $O(r^4)=O(x^6)$. The matrices $L_P,R_P$ are bounded
polynomials, $s-m=O(x)$ and $\ell_0=O(x)$. Also
\[
 |f|\geq2mr-2\lambda=2r(m-r)\geq c_1xr,
 \qquad |d|\geq c_2x^4
\]
for positive constants $c_1,c_2$ and sufficiently small $x>0$.
Indeed $m/x\to54/25$ while $r/x\to0$.
It follows from~\eqref{eq:beta-s} that $B_s=O(x^{-4}/r)$ uniformly
on the circle. Including its length $2\pi r$ and the factor $1/\tau$,
\begin{equation}\label{eq:upper-beta}
 \mathfrak b(U)=O(x^2)\longrightarrow0.
\end{equation}

\subsection{All singular and constant terms of the outer correction cancel}
The literal matrix $C$ gives
\[
 C^{-1}(1,0,0,0,0,0)^{\tr}=v_0=(0,0,0,5/72,0,0)^{\tr},
\]
\[
 C^{-1}(1,1,1,1,1,1)^{\tr}
   =v_1=(-1,-1,-1,77/24,77/24,77/24)^{\tr}.
\]
Since $a_1^2\rho=9/40000$, the normally convergent series give
\begin{equation}\label{eq:state-expansion}
 \vv(\rho x^3)=v_0+\frac9{40000}x^3v_1+O(x^6).
\end{equation}
Expand the full rational row as $\xi(x)=\sum_{k=-4}^{\infty}\xi_kx^k$
near zero. Exact substitution of its printed coefficients yields
\begin{align}
 \xi_{-4}v_0=\xi_{-3}v_0=\xi_{-2}v_0&=0,\label{eq:outer-cancel-a}\\
 \xi_{-1}v_0=-\frac{105800}{459459},\quad
 \frac9{40000}\xi_{-4}v_1&=\frac{105800}{459459},\label{eq:outer-cancel-b}\\
 \xi_0v_0=\frac{8464}{21879},\quad
 \frac9{40000}\xi_{-3}v_1&=-\frac{8464}{21879}.
 \label{eq:outer-cancel-c}
\end{align}
Both $(99x-50)^{-1}$ and $(99x-50)^{-2}$ contribute to $\xi_0$.
Equations~\eqref{eq:outer-cancel-a}--\eqref{eq:outer-cancel-c} cancel
all negative powers and the constant term of $\xi\vv$.
Multiplying the $O(x^6)$ remainder in~\eqref{eq:state-expansion} by
$\xi=O(x^{-4})$ gives $O(x^2)$. Thus $\xi\vv=O(x)$ and
\begin{equation}\label{eq:upper-full}
 \mathcal E(1-)=0.
\end{equation}
At the other endpoint $x\to1$, the row $\xi$ is finite and the state is
analytic at $\lambda=\rho$. Consequently
\begin{equation}\label{eq:lower-outer}
 U\xi(x)\vv(\lambda)\longrightarrow0\qquad(U\to0^+).
\end{equation}

\section{The lower endpoint is a full positive residue}
Choose once and for all $\rho<r_0<u$; for example $r_0=1/10$.
For sufficiently small $U>0$, both roots of $f$ lie outside $|s|=r_0$,
whereas the balanced circle encloses $\sigma_+$ alone. On the fixed
inner circle $B_s\to0$ uniformly: it has the explicit factor $U$, its
remaining denominator is nonzero at $U=0$, and both native arguments
stay in compact subsets of the unit disk. Hence its normalized integral
tends to zero.

The positively oriented annulus between the two circles contains exactly
the pole $\sigma_+$. The residue theorem gives their difference as the
full residue of the scalar integrand in~\eqref{eq:beta-flux}, because
$2\pi i/\tau=1$. Algebraically,
\begin{equation}\label{eq:inside-residue}
 \Res_{s=\sigma_+}B_s
 =\frac{U L_P+\sqrt M R_P}{d}
 \longrightarrow S_*:=\frac{4R_P(1)}{d(1)}.
\end{equation}
No pole lies on either contour in this application. Taking the limit of
the actual inside residue gives
\begin{equation}\label{eq:lower-beta}
 \mathfrak b(0+)=Y(u)^{\tr}T(u)S_*T(u)^{\tr}Y(u).
\end{equation}
The polynomial matrices yield the small exact product
\begin{equation}\label{eq:residue-matrix}
 T(u)S_*T(u)^{\tr}=
 \begin{pmatrix}
 -1/750&1/150&98/1125\\
 1/150&98/1125&0\\98/1125&0&0
 \end{pmatrix}.
\end{equation}
Every quantity in this computation is rational, except the subsequently
evaluated native solution values. The sign is fixed by the annulus
orientation and the selected inside root.

The native Gauss equation and symmetric-square identity proved in the
next section imply, at $q=u$,
\[
 F=g^2,\qquad \theta F=2g\theta g,\qquad
 \theta^2 F=2(\theta g)^2+2g\theta^2g,
\]
with
$\theta^2g=(27/196)\theta g+(15/1568)g$ at $u$.
Substitution in~\eqref{eq:residue-matrix} gives
\begin{equation}\label{eq:residue-square}
 \mathfrak b(0+)=\frac{\{3F(u)+28\theta F(u)\}^2}{4500}.
\end{equation}
For an explicit check, the coefficients of $g^4,g^3\theta g,
g^2(\theta g)^2$ on both sides are respectively
$9/4500,336/4500,3136/4500$.

% Claim: the actual native identity 3 F(27/125)+28 theta F(27/125)=5 sqrt(5)/pi,
% including the signed algebraic Gauss transformation for every real h>8.
% Scope: the literal convergent hypergeometric functions defined below and
% their ordinary derivatives on the stated real intervals.
% Evidence: [ARG], the complete human-checkable analytic derivation below.
% Proof status: Complete for this scalar analytic lemma.
% Open premises: None within this lemma.
% Review: Not reviewed in this fragment; the assembled document receives
% the parent's single fresh adversarial mathematical review.
% Human acceptance: Pending.

\section{A native Gauss certificate for the \texorpdfstring{$1/\pi$}{1/pi} value}
\label{sec:native-gauss-certificate}

Put
\begin{equation}\label{eq:gauss-definitions}
\begin{aligned}
 H(p)&=\sum_{n\geq0}\frac{(1/6)_n(5/6)_n}{(n!)^2}p^n,\\
 F(q)&=\sum_{n\geq0}
       \frac{(1/6)_n(1/2)_n(5/6)_n}{(n!)^3}q^n,
       \qquad \theta=q\frac{d}{dq},\quad \tau=2\pi i.
\end{aligned}
\end{equation}
Here $(a)_0=1$ and $(a)_n=a(a+1)\cdots(a+n-1)$ for $n\geq1$.
Both series have radius of convergence one, by their consecutive
coefficient ratios.  Thus they and their termwise derivatives define
analytic functions on the open unit disc.  For real $0<q<1$, define
\begin{equation}\label{eq:gauss-p-g-Z}
 p(q)=\frac{1-\sqrt{1-q}}2,
 \qquad g(q)=H(p(q)),\qquad Z(q)=H(1-p(q)).
\end{equation}
Every square root and fractional power of a positive real number below
has its positive real value.  A prime on $H$ always means its ordinary
derivative with respect to its argument.

\subsection{The normalized symmetric square}

The coefficient recurrence of $H$ gives the actual Gauss equation
\begin{equation}\label{eq:gauss-ordinary-equation}
 p(1-p)H''(p)+(1-2p)H'(p)-\frac5{36}H(p)=0,
 \qquad H(0)=1.
\end{equation}
Indeed its Euler form is
$\{p\partial_p\}^2H-p\{p\partial_p+1/6\}
\{p\partial_p+5/6\}H=0$, whose coefficient of $p^n$ is
$n^2h_n-(n-5/6)(n-1/6)h_{n-1}$ for $n\geq1$.
The equation is invariant under $p\mapsto1-p$, so $H(1-p)$ also
satisfies it wherever $0<p<1$.

Write $s=\sqrt{1-q}$.  The identities
\[
 p(1-p)=q/4,\qquad 1-2p=s,\qquad
 \frac{dp}{dq}=\frac1{4s},\qquad
 \frac{d^2p}{dq^2}=\frac1{8s^3}
\]
and the chain rule give, for either $y=g$ or $y=Z$,
\begin{equation}\label{eq:gauss-native}
 \theta^2y=\frac{q}{2(1-q)}\theta y+
                 \frac{5q}{144(1-q)}y.
\end{equation}
For clarity, the corresponding ordinary derivative equation is
\[
 y''=\left(-\frac1q+\frac1{2(1-q)}\right)y'
                   +\frac5{144q(1-q)}y,
\]
where these primes mean $q$ derivatives.

Define the two differential expressions
\[
\begin{aligned}
 \mathcal E_2y&=144(1-q)\theta^2y-72q\theta y-5qy,\\
 \mathcal E_3f&=72(1-q)\theta^3f
                      -q(108\theta^2f+46\theta f+5f).
\end{aligned}
\]
The product rule, without any transformation formula as an input, gives
\begin{equation}\label{eq:gauss-symmetric-square-identity}
 \mathcal E_3(y^2)
       =y\theta(\mathcal E_2y)+3(\theta y)\mathcal E_2y.
\end{equation}
For example, the expansion of either side is
\[
144(1-q)y\theta^3y+432(1-q)(\theta y)(\theta^2y)
-216qy\theta^2y-216q(\theta y)^2-92qy\theta y-5qy^2.
\]
Consequently $g^2$ is annihilated by $\mathcal E_3$.
The series $F$ is annihilated by the same expression, since
\begin{equation}\label{eq:gauss-native-recurrence}
 n^3 f_n=(n-5/6)(n-1/2)(n-1/6)f_{n-1}
 \qquad(n\geq1).
\end{equation}
This recurrence uniquely determines every analytic solution from its
constant coefficient: the multiplier $n^3$ is nonzero for $n\geq1$.
The function $g$ is analytic near zero because $p(0)=0$ and the square
root in \eqref{eq:gauss-p-g-Z} is analytic there.  Since $g(0)^2=F(0)=1$,
the recurrence proves equality near zero.  Both sides are real analytic
on the connected interval $(0,1)$, so
\begin{equation}\label{eq:symmetric-square}
                         F(q)=g(q)^2\qquad(0<q<1).
\end{equation}
This last continuation can equivalently be obtained by uniqueness for
the nonsingular third-order ordinary equation on each compact
subinterval of $(0,1)$.

\subsection{A signed algebraic transformation}

For every real $h>8$, set
\begin{equation}\label{eq:gauss-h-parameters}
\begin{aligned}
 v_1(h)&=\frac{(h-8)\sqrt{h+64}}{(h+16)^{3/2}},&
 p_1(h)&=\frac{1-v_1(h)}2,&
 q_1(h)&=\frac{1728h}{(h+16)^3},\\
 v_2(h)&=\frac{(h-512)\sqrt{h+64}}{(h+256)^{3/2}},&
 p_2(h)&=\frac{1-v_2(h)}2,&
 q_2(h)&=\frac{1728h^2}{(h+256)^3},\\
 L(h)&=\left(\frac{h+256}{h+16}\right)^{1/4}.&&&
\end{aligned}
\end{equation}
We will prove the identity of actual functions
\begin{equation}\label{eq:gauss-signed-transformation}
                         H(p_2(h))=L(h)H(p_1(h))
                         \qquad(h>8).
\end{equation}
The numerator $h-512$ is signed throughout this assertion.

Here is a direct pullback and normalization proof, including the
rational identities needed to check its gauge.  Put $z=1/h$ and
\[
 S=1+64z,\qquad d_1=1+16z,\qquad d_2=1+256z,
 \qquad \Delta=z\frac{d}{dz}.
\]
The functions $\widetilde p_i(z)=p_i(1/z)$ extend analytically to zero
by the explicit formulas
\begin{equation}\label{eq:gauss-z-parameters}
 \widetilde p_1=\frac{1-(1-8z)\sqrt S/d_1^{3/2}}2,
 \qquad
 \widetilde p_2=\frac{1-(1-512z)\sqrt S/d_2^{3/2}}2.
\end{equation}
The elementary polynomial identities
\begin{equation}\label{eq:gauss-z-discriminants}
 d_1^3-(1-8z)^2S=1728z^2,
 \qquad d_2^3-(1-512z)^2S=1728z
\end{equation}
imply
\[
 \widetilde p_1(1-\widetilde p_1)=\frac{432z^2}{d_1^3},
 \qquad
 \widetilde p_2(1-\widetilde p_2)=\frac{432z}{d_2^3}.
\]
In particular $0<\widetilde p_i<1$ when $0<z<1/8$; the first one is
also less than $1/2$ because $1-8z>0$.

Differentiating the literal signed expressions gives
\begin{equation}\label{eq:gauss-corrected-z-derivatives}
 \frac{d\widetilde p_1}{dz}=\frac{864z}{\sqrt S\,d_1^{5/2}},
 \qquad
 \frac{d\widetilde p_2}{dz}=\frac{432}{\sqrt S\,d_2^{5/2}}.
\end{equation}
For an explicit check of these derivative numerators,
\[
\begin{aligned}
 -8Sd_1+(1-8z)(32d_1-24S)&=-1728z,\\
 -512Sd_2+(1-512z)(32d_2-384S)&=-864.
\end{aligned}
\]
Thus there is no extra factor $1-8z$ in either derivative.  Their
logarithmic derivatives on $0<z<1/8$ are
\begin{equation}\label{eq:gauss-z-log-derivatives}
 \frac{\widetilde p_1''}{\widetilde p_1'}
     =\frac1z-\frac{32}{S}-\frac{40}{d_1},
 \qquad
 \frac{\widetilde p_2''}{\widetilde p_2'}
     =-\frac{32}{S}-\frac{640}{d_2},
\end{equation}
where the derivatives in this display are with respect to $z$.

For any such pullback $Y(z)=H(\widetilde p(z))$, the chain rule in
\eqref{eq:gauss-ordinary-equation} yields
\[
 \Delta^2Y+B\Delta Y+CY=0,
\]
with the explicit coefficients
\begin{equation}\label{eq:gauss-general-pullback}
 B=\frac{z\widetilde p'(1-2\widetilde p)}
               {\widetilde p(1-\widetilde p)}
               -1-\frac{z\widetilde p''}{\widetilde p'},
 \qquad
 C=-\frac5{36}\frac{z^2(\widetilde p')^2}
                         {\widetilde p(1-\widetilde p)}.
\end{equation}
All divisions here are legitimate by
\eqref{eq:gauss-z-discriminants} and
\eqref{eq:gauss-corrected-z-derivatives}.  Substitution gives
\begin{equation}\label{eq:gauss-corrected-z-coefficients}
\begin{aligned}
 B_1&=z\left(-\frac8{d_1}+\frac{32}{S}\right),&
 C_1&=-\frac{240z^2}{Sd_1^2},\\
 B_2&=z\left(-\frac{128}{d_2}+\frac{32}{S}\right),&
 C_2&=-\frac{60z}{Sd_2^2}.
\end{aligned}
\end{equation}
For instance the first fractions in the expressions for $B_1,B_2$
in \eqref{eq:gauss-general-pullback} are respectively
$2(1-8z)/d_1$ and $(1-512z)/d_2$; together with
\eqref{eq:gauss-z-log-derivatives} these produce exactly the displayed
$B_i$.

Let
\[
 \mathcal L(z)=\left(\frac{d_2}{d_1}\right)^{1/4},
 \qquad \mu(z)=\frac{\mathcal L'(z)}{\mathcal L(z)}
                  =\frac{60}{d_1d_2}.
\]
Its ordinary $z$ derivative is
\[
 \mu'=-\frac{60(16d_2+256d_1)}{d_1^2d_2^2}
      =-\frac{60(272+8192z)}{d_1^2d_2^2}.
\]
The gauge identities are
\begin{equation}\label{eq:gauss-z-gauge-identities}
 B_1=B_2+2z\mu,
 \qquad
 C_1=C_2+z^2(\mu'+\mu^2)+z(1+B_2)\mu.
\end{equation}
The first follows from $-8d_2+128d_1=120$.  To verify the second
without an implicit computational certificate, multiply it by
$Sd_1^2d_2^2$ and use this polynomial identity:
\begin{equation}\label{eq:gauss-z-gauge-polynomial}
\begin{aligned}
 -240z^2d_2^2+60zd_1^2
  ={}&Sz^2(3600-960d_2-15360d_1)
       +60zSd_1d_2\\
     &{}-7680z^2Sd_1+1920z^2d_1d_2.
\end{aligned}
\end{equation}
It is an identity after $S=1+64z$, $d_1=1+16z$ and $d_2=1+256z$
are inserted.

We also give the ordinary $h$ equations explicitly.  Since
$\Delta=-h\,d/dh$ on a function pulled back by $z=1/h$, the corrected
coefficients \eqref{eq:gauss-corrected-z-coefficients} imply
\begin{equation}\label{eq:gauss-h-equations}
 Y_i''+a_iY_i'+b_iY_i=0,\qquad Y_i(h)=H(p_i(h)),
\end{equation}
where these primes are $h$ derivatives and
\begin{equation}\label{eq:gauss-h-coefficients}
\begin{aligned}
 a_1&=\frac1h\left(1+\frac8{h+16}-\frac{32}{h+64}\right),&
 b_1&=-\frac{240}{h(h+64)(h+16)^2},\\
 a_2&=\frac1h\left(1+\frac{128}{h+256}-\frac{32}{h+64}\right),&
 b_2&=-\frac{60}{(h+64)(h+256)^2}.
\end{aligned}
\end{equation}
Indeed $a_i=(1-B_i(1/h))/h$ and $b_i=C_i(1/h)/h^2$.
Writing
\begin{equation}\label{eq:gauss-h-lambda}
 \mu_h(h)=\frac{L'(h)}{L(h)}
      =-\frac{60}{(h+16)(h+256)},
 \qquad
 \mu_h'(h)=\frac{60(2h+272)}{(h+16)^2(h+256)^2},
\end{equation}
the same identities become
\begin{equation}\label{eq:gauss-h-gauge-identities}
 a_1=a_2+2\mu_h,
 \qquad b_1=b_2+\mu_h'+\mu_h^2+a_2\mu_h.
\end{equation}
For the second conversion, explicitly
$\mu_h=-\mu(1/h)/h^2$ and
\[
 \mu_h'+\mu_h^2+a_2\mu_h
 =\frac{\mu'(1/h)+\mu(1/h)^2}{h^4}
   +\frac{(1+B_2(1/h))\mu(1/h)}{h^3}
 =\frac{C_1(1/h)-C_2(1/h)}{h^2}.
\]
Thus the product rule proves that $LY_1$ satisfies the second
ordinary equation:
\[
 (LY_1)''+a_2(LY_1)'+b_2LY_1
 =L\{Y_1''+(2\mu_h+a_2)Y_1'
          +(\mu_h'+\mu_h^2+a_2\mu_h+b_2)Y_1\}=0.
\]

It remains to fix the actual solution, rather than only its equation.
Both $H(\widetilde p_2(z))$ and
$\mathcal L(z)H(\widetilde p_1(z))$ are analytic near zero, take the
value one there, and satisfy
$\Delta^2Y+B_2\Delta Y+C_2Y=0$; the displayed formulas give
$\widetilde p_1(0)=\widetilde p_2(0)=0$ and $\mathcal L(0)=1$.
The rational functions $B_2,C_2$
are analytic and vanish at zero.  If their two solutions differed,
the first nonzero term of their analytic difference would be
$cz^m$ with $m\geq1$.  The coefficient of $z^m$ in its equation
would then be $m^2c$, because the other two terms have higher order.
This contradicts $c\ne0$.  Hence the solutions agree near $z=0$.
On $h>8$ the coefficients \eqref{eq:gauss-h-coefficients} are
continuous and the equation is nonsingular.  Ordinary initial-value
uniqueness continues the equality from sufficiently large $h$ to
the entire connected interval $h>8$, proving
\eqref{eq:gauss-signed-transformation}.  This uniqueness step requires
only the usual local uniqueness for a linear first-order system:
on a compact interval its bounded coefficient matrix gives uniqueness
by the integral equation and Gronwall's inequality.

All functions in \eqref{eq:gauss-h-parameters} are smooth at $h=512$.
In particular $p_2(512)=1/2$ and $H$ is analytic there; the proof makes
no change of branch at that point.  Differentiating the actual identity
now gives the fully explicit ordinary derivative identity
\begin{equation}\label{eq:gauss-actual-h-derivative}
 -\frac{432h\,H'(p_2(h))}{\sqrt{h+64}(h+256)^{5/2}}
 =L(h)\left\{
     -\frac{60H(p_1(h))}{(h+16)(h+256)}
     -\frac{864H'(p_1(h))}{\sqrt{h+64}(h+16)^{5/2}}
          \right\}.
\end{equation}
Here the two chain derivatives are exactly
$p_1'(h)=-864/[\sqrt{h+64}(h+16)^{5/2}]$ and
$p_2'(h)=-432h/[\sqrt{h+64}(h+256)^{5/2}]$.

\subsection{The Euler endpoint and the normalized Wronskian}

We derive the normalization from an integral, including its
differentiated endpoint.  For $x,y>0$ write
$B(x,y)=\int_0^1t^{x-1}(1-t)^{y-1}\,dt$.
The change of variable $t=r/(1+r)$ and then $r=w^6$ gives
\[
 B(5/6,1/6)=\int_0^\infty\frac{r^{-1/6}}{1+r}\,dr
                         =6\int_0^\infty\frac{w^4}{1+w^6}\,dw.
\]
Here is an elementary residue evaluation of the last integral.  The
upper half-plane poles of $w^4/(1+w^6)$ are
$\exp(i\pi/6),i,\exp(5i\pi/6)$, and the residue at a pole $\zeta$
is $1/(6\zeta)$.  Their sum is
\[
 \frac16\{e^{-i\pi/6}+e^{-i\pi/2}+e^{-5i\pi/6}\}
                         =-\frac i3.
\]
For a semicircle of radius $R\geq2$ the absolute value of the
integrand is at most $2/R^2$, so its arc integral tends to zero.
The residue theorem therefore gives
$\int_{-\infty}^{\infty}w^4/(1+w^6)\,dw=2\pi/3$.
The integrand is even, and consequently
\begin{equation}\label{eq:gauss-beta-normalization}
                              B(5/6,1/6)=2\pi.
\end{equation}

For $0\leq x<1$, binomial expansion in the positive integral yields
\begin{equation}\label{eq:gauss-euler-integral}
 H(x)=\frac1{2\pi}\int_0^1
       t^{-1/6}(1-t)^{-5/6}(1-xt)^{-1/6}\,dt.
\end{equation}
To see the coefficients without assuming an Euler formula, integration
by parts gives $B(a+1,b)=aB(a,b)/(a+b)$ for $a,b>0$.
Indeed the integral of $(t^a(1-t)^b)'$ is zero, and
$B(a,b)=B(a+1,b)+B(a,b+1)$.
It follows that
$B(5/6+n,1/6)/B(5/6,1/6)=(5/6)_n/n!$.
Inserting the binomial coefficients $(1/6)_n/n!$ in the integral
therefore produces exactly \eqref{eq:gauss-definitions}; monotone
convergence justifies the interchange for $0\leq x<1$.

On every compact interval $0\leq x\leq r<1$, differentiation under
this integral is justified by the integrable majorant
\[
 \frac{(1-r)^{-7/6}}{12\pi}
                       t^{5/6}(1-t)^{-5/6}.
\]
Thus the derivative of the actual function is
\begin{equation}\label{eq:gauss-differentiated-euler}
 H'(x)=\frac1{12\pi}\int_0^1
       t^{5/6}(1-t)^{-5/6}(1-xt)^{-7/6}\,dt.
\end{equation}
Take $x=1-\delta$ with $0<\delta\leq1/2$, first set $r_0=1-t$
and then $r_0=\delta v$.  This gives the exact equation
\begin{equation}\label{eq:gauss-scaled-derivative-integral}
 \delta H'(1-\delta)=\frac1{12\pi}\int_0^{1/\delta}
    (1-\delta v)^{5/6}v^{-5/6}
                  \{1+(1-\delta)v\}^{-7/6}\,dv.
\end{equation}
Extend this integrand by zero for $v>1/\delta$.  It is bounded by
$v^{-5/6}(1+v/2)^{-7/6}$, which is integrable: it is $O(v^{-5/6})$
at zero and $O(v^{-2})$ at infinity.  Dominated convergence and the
substitution $w=v/(1+v)$ now give
\[
 \int_0^\infty v^{-5/6}(1+v)^{-7/6}\,dv
                      =\int_0^1w^{-5/6}\,dw=6.
\]
Consequently
\begin{equation}\label{eq:gauss-differentiated-endpoint}
                       \lim_{\delta\downarrow0}
                       \delta H'(1-\delta)=\frac1{2\pi}.
\end{equation}

The undifferentiated integral has a sufficient elementary bound of its
own: $1-(1-\delta)t\geq\delta$ for $0\leq t\leq1$.  By
\eqref{eq:gauss-beta-normalization} and
\eqref{eq:gauss-euler-integral},
\begin{equation}\label{eq:gauss-endpoint-bound}
 0<H(1-\delta)\leq\delta^{-1/6},\qquad
 0<\delta H(1-\delta)\leq\delta^{5/6}\longrightarrow0
 \qquad(\delta\downarrow0).
\end{equation}

For $0<q<1$ define the Euler Wronskian
\[
                \mathcal W(q)=g(q)\theta Z(q)-(\theta g(q))Z(q).
\]
Equation \eqref{eq:gauss-native} gives
\[
 \theta\mathcal W=\frac{q}{2(1-q)}\mathcal W,
 \qquad \frac{d}{dq}\{\sqrt{1-q}\,\mathcal W(q)\}=0.
\]
Its constant is fixed by \eqref{eq:gauss-differentiated-endpoint}:
as $q\downarrow0$, $p(q)=q/4+O(q^2)$ and
\[
 \frac{q\,p'(q)}{p(q)}=\frac{1+\sqrt{1-q}}{2\sqrt{1-q}}\longrightarrow1,
 \qquad
 \theta Z(q)=-q\,p'(q)H'(1-p(q))\longrightarrow-\frac1{2\pi}.
\]
Also $g(q)\to1$ and $\theta g(q)=q\,p'(q)H'(p(q))=O(p(q))$:
$H'$ is bounded near zero and $q/p(q)\to4$.
The independent bound \eqref{eq:gauss-endpoint-bound} therefore gives
$(\theta g(q))Z(q)=O(p(q)^{5/6})\to0$.
We have proved
\begin{equation}\label{eq:gauss-normalized-wronskian}
 g(q)\theta Z(q)-(\theta g(q))Z(q)
                         =-\frac1{2\pi\sqrt{1-q}}.
\end{equation}
In particular the derivative normalization comes directly from the
differentiated integral and its integrable bound, not from
differentiating a remainder in a bare logarithmic expansion.

\subsection{The actual contact at \texorpdfstring{$h=64$}{h=64}}

Put $u=27/125$ and $s_u=\sqrt{1-u}=7\sqrt{10}/25$.
Substitution in \eqref{eq:gauss-h-parameters} gives
\begin{equation}\label{eq:gauss-contact-parameters}
\begin{aligned}
 q_1(64)&=q_2(64)=u,&L(64)&=\sqrt2,\\
 v_1(64)&=s_u,&v_2(64)&=-s_u,\\
 p_1(64)&=p(u),&p_2(64)&=1-p(u).
\end{aligned}
\end{equation}
Therefore the actual transformation already proves
\begin{equation}\label{eq:gauss-contact-value}
                               Z(u)=\sqrt2\,g(u).
\end{equation}
On a neighborhood of $h=64$ contained in $(8,512)$ the positive-root
definition in \eqref{eq:gauss-p-g-Z} satisfies
$p_1(h)=p(q_1(h))$ and $p_2(h)=1-p(q_2(h))$.
This follows from $4p_i(1-p_i)=q_i$, $v_1>0$ and $v_2<0$,
which are also consequences of \eqref{eq:gauss-z-discriminants}.
We may consequently differentiate the actual equality
$Z(q_2(h))=L(h)g(q_1(h))$ on that neighborhood.
The three logarithmic derivatives are
\[
 \frac{q_1'}{q_1}=\frac1h-\frac3{h+16},\qquad
 \frac{q_2'}{q_2}=\frac2h-\frac3{h+256},\qquad
 \frac{L'}L=-\frac{60}{(h+16)(h+256)},
\]
so at $h=64$ their values are $-7/320$, $7/320$ and $-3/1280$.
The differentiated identity is thus
\[
 \frac7{320}\theta Z(u)
       =\sqrt2\left\{-\frac3{1280}g(u)
                                  -\frac7{320}\theta g(u)\right\},
\]
or, equivalently,
\begin{equation}\label{eq:gauss-contact-theta}
             \theta Z(u)=-\sqrt2\,\theta g(u)
                                      -\frac{3\sqrt2}{28}g(u).
\end{equation}
This relation uses the ordinary derivative of an identity on an open
interval, with all moving parameters and the derivative of $L$ retained.

Finally insert \eqref{eq:gauss-contact-value} and
\eqref{eq:gauss-contact-theta} into
\eqref{eq:gauss-normalized-wronskian}, and use
$F=g^2$ and $\theta F=2g\theta g$ from
\eqref{eq:symmetric-square}.  It follows that
\[
 -\sqrt2\left\{\theta F(u)+\frac3{28}F(u)\right\}
                         =-\frac1{2\pi s_u}.
\]
Since $s_u=7\sqrt{10}/25$, this is the required native certificate
\begin{equation}\label{eq:native-value}
 \boxed{\displaystyle
             3F\left(\frac{27}{125}\right)
                +28\theta F\left(\frac{27}{125}\right)
                           =\frac{5\sqrt5}{\pi}.}
\end{equation}

\section{Completion of the series evaluation}
Equation~\eqref{eq:native-value} substituted in~\eqref{eq:residue-square}
gives
\[
 \mathfrak b(0+)=\frac{(5\sqrt5/\pi)^2}{4500}
              =\frac{1}{36\pi^2}.
\]
Together with~\eqref{eq:lower-outer}, this gives
$\mathcal E(0+)=1/(36\pi^2)$. The upper limit is
\eqref{eq:upper-full}, and the escape singularity has been removed.
On every closed subinterval of $(0,1)$ the extended primitive satisfies
\eqref{eq:primitive-derivative}. The loaded radial integrand is continuous
and bounded on $[0,1]$, by its uniformly convergent polynomial-geometric
series. Applying the fundamental theorem of calculus first away from
the endpoints and then taking their limits yields
\[
 \int_0^1Q(\theta_\lambda)h\bigl(\rho(1-U^2)^3\bigr)\,\dd U
   =3375\{\mathcal E(0+)-\mathcal E(1-)\}
   =\frac{3375}{36\pi^2}=\frac{375}{4\pi^2}.
\]
Finally~\eqref{eq:radial} identifies this entire integral with the
literal factorial series in Theorem~\ref{thm:main}, proving the theorem.

\appendix
\section{A finite rational certificate for the source row}
\label{app:certificates}

This appendix supplies the rational identities used in the contour
comparison.  In particular, it includes an actual primitive, rather than a
rank computation or a numerical fit.  Everything in the certificate is
defined below over rational function fields.  The contour's analytic
boundary and normalization are the statements proved in the main text;
finite algebra by itself would not establish those statements.

Write
\[
 \rho=\frac{729}{15625},\qquad x=1-U^2,\qquad
 \lambda=\rho x^3,\qquad b=\frac{\lambda}{s},\qquad
 \tau=2\pi i.
\]
All coefficient rows act on solution columns.  Thus their covariant
differential has a \emph{plus} sign.  Tensor coordinates are ordered
\((00,01,02,10,11,12,20,21,22)\).

\subsection{The complete source coefficient}

The ordinary elliptic symmetric-square state is
\((y^2,2yy',y'^2)^{\mathsf T}\); its third entry is not the second
derivative of its first entry.  Its connection and Euler gauge are
\[
 \mathsf A(q)=\begin{pmatrix}0&1&0\\2c(q)&-a(q)&2\\0&c(q)&-2a(q)\end{pmatrix},
 \quad a(q)=\frac{1-3q/2}{q(1-q)},\quad
 c(q)=\frac{5}{144q(1-q)},
\]
\[
 \mathsf G(q)=\begin{pmatrix}
 1&0&0\\0&q&0\\
 \dfrac{5q}{72(1-q)}&\dfrac{q^2}{2(1-q)}&2q^2
 \end{pmatrix},\qquad
 N(q)=\begin{pmatrix}0&1&0\\0&0&1\\
 \dfrac{5q}{72(1-q)}&\dfrac{23q}{36(1-q)}&\dfrac{3q}{2(1-q)}
 \end{pmatrix}.
\]
They satisfy \(\mathsf G'+\mathsf G\mathsf A=(N/q)\mathsf G\).  The solution-state form and its useful
product with the gauge are
\[
 J_C(q)=\begin{pmatrix}
 -5q/36&-q/2&1-q\\-q/2&q-1&0\\1-q&0&0
 \end{pmatrix}.
\]
\[
 T(q)=J_C(q)\mathsf G(q)=\begin{pmatrix}
 -5q/72&0&2q^2(1-q)\\-q/2&q(q-1)&0\\1-q&0&0
 \end{pmatrix}.
\]

Here are the coefficient polynomials of the literal rational source.
Put
\begin{align*}
 d&=x^4(25-9x)(99x-50)^2,&
 K_2&=81x^2-198x+80,\\
 K_3&=243x^3-1053x^2+1266x-400,&
 K_6&=59049x^6-511758x^5+1724085x^4\\
 &&&\quad-2860596x^3+2441988x^2-1011200x+160000,\\
 k&=\frac{1953125}{209952},&j&=\frac{244140625}{944784},\\
 v&=\frac{3906250}{6561},&c_*&=\frac{30517578125}{8503056}.
\end{align*}
Then
\[
 L_P=\begin{pmatrix}
 0&0&-k(9x-25)K_2K_3\\
 0&0&j(9x-25)K_2\\
 k(9x-25)K_2K_3&-j(9x-25)K_2&0
 \end{pmatrix}.
\]
\[
 R_P=\begin{pmatrix}
 0&0&kK_6\\0&vx(99x-50)&-jK_3\\kK_6&-jK_3&c_*
 \end{pmatrix}.
\]
For compact notation put \(L=U L_P/d\) and \(R=R_P/d\).  Define
\[
 m=\frac{27x(80-99x+27x^2)}{1000},\quad
 \ell_0=\frac{27x(9x-16)}{1000},\quad
 \ell=U\ell_0,\quad M=25-9x,
\]
\[
 z=s-m,\qquad f=z^2-\ell^2M=s^2-2ms+\lambda.
\]
The entire rational coefficient one-form is
\begin{equation}\label{eq:appendix-beta}
 \beta_N=\frac{L\,\dd f+R(2\ell M\,\dd z-2zM\,\dd\ell-z\ell\,\dd M)}{f}
           =T_U\,\dd U+U T_s\,\dd s,
\end{equation}
where, with primes denoting \(\partial_x\) at fixed \(s\),
\begin{align}
 T_s&=\frac{2L_P(s-m)+2\ell_0M R_P}{df},\label{eq:appendix-Ts}\\
 T_U&=\frac{-2(1-x)L_P f_x
 -2R_P\{-2\ell_0(1-x)Mm'+(s-m)[(\ell_0-2(1-x)\ell_0')M+9(1-x)\ell_0]\}}{df}.
 \label{eq:appendix-TU}
\end{align}
Thus the correction has not been omitted or absorbed into an unspecified
boundary term.  The relation to the main coefficients is \(B_s=UT_s\) and \(B_U=T_U\).
Taking their covariant derivatives gives
\begin{align}
 \mathsf K={}&-T_s+2(1-x)\partial_xT_s
 -\frac{6(1-x)b}{x}T_s \mathsf A(b)^{\mathsf T}\notag\\
 &+\partial_sT_U-\mathsf A(s)T_U+\frac bsT_U \mathsf A(b)^{\mathsf T},
 \label{eq:appendix-alpha}\\
 \mathfrak q&=\operatorname{flatten}\bigl(T(s)\mathsf K T(b)^{\mathsf T}\bigr).
 \label{eq:appendix-Q}
\end{align}
In particular, the source row against the native product
Euler state is \(\mathfrak q\,\dd U\wedge\dd s\).
Thus \(\mathfrak q=\operatorname{flatten}(\mathsf S)\) in
the notation of the main text. Formula
\eqref{eq:appendix-alpha} is simply
\(-\nabla_U B_s+\nabla_s B_U\), with \(x_U=-2U\) and
\(b_U=-6Ub/x\); these identities fix both its sign and its chain-rule
terms.

\subsection{Angular normal form and its six-state action}

Let
\[
 J=\begin{pmatrix}0&1&0\\0&0&1\\0&0&0\end{pmatrix},\qquad
 \ww=\begin{pmatrix}5/72&23/36&3/2\end{pmatrix},\qquad E=\e_2^{\tr}\ww,
\]
Here \(\e_2=(0,0,1)\) is the standard row, as in the main text.  Over \(\mathbb Q(\lambda)(s)\),
\[
 \mathcal M_s=\frac{N(s)\otimes I-I\otimes N(\lambda/s)}s
 =\frac{R_0}s+\frac{R_1}{s-1}+\frac{R_\lambda}{s-\lambda},
\]
\[
 R_0=J\otimes I-I\otimes(J-E),\qquad
 R_1=-E\otimes I,\qquad R_\lambda=-I\otimes E.
\]
For a six-row \(v=(a_0,a_1,a_2,b_0,b_1,b_2)=(a,b_*)\), define
\begin{equation}\label{eq:appendix-NF}
 \operatorname{NF}(v)=\frac{\ww\otimes a}{s-1}
                         +\frac{b_*\otimes \ww}{s-\lambda}.
\end{equation}
The six standard coordinate rows under this formula are the six angular
basis representatives.  Their induced \(D_\lambda\) action is
\begin{equation}\label{eq:appendix-B6}
 B_6(\lambda)=\begin{pmatrix}
 J+\dfrac{\lambda E}{1-\lambda}&-\dfrac{\lambda E}{1-\lambda}\\[2pt]
 -\dfrac{E}{1-\lambda}&J+\dfrac{\lambda E}{1-\lambda}
 \end{pmatrix}.
\end{equation}
This assertion includes an explicit angular primitive.  For a constant
row \(v=(a,b_*)\), put
\(p_v=-s(b_*\otimes \ww)/(s-\lambda)\).  Entrywise,
\begin{equation}\label{eq:appendix-six-action}
 \lambda\partial_\lambda\operatorname{NF}(v)
 +\operatorname{NF}(v)(I\otimes N(\lambda/s))
 -\operatorname{NF}(vB_6)
 =\partial_sp_v+p_v\mathcal M_s.
\end{equation}
For the first three basis rows \(p_v=0\).

The coefficient one-form for the anchor is \(\omega_0=e_{00}\,\dd s/s\),
whose normalized contour integral is the scalar function \(h\).  If
\(\eta_0=e_{00}R_0^{-1}\), its exact angular identity is
\(e_{00}/s=\operatorname{NF}(c_0)+\eta_0\mathcal M_s\).  The coordinate rows of
\(\omega_0,D_\lambda\omega_0,\ldots,D_\lambda^5\omega_0\) form the following full matrix:
\begin{equation}\label{eq:appendix-C}
 C=\begin{pmatrix}
 49456/125&153648/125&113472/125&72/5&-3312/25&113472/125\\
 -1576/25&-4608/25&-3312/25&0&72/5&-3312/25\\
 46/5&108/5&72/5&0&0&72/5\\
 -1&0&0&0&0&0\\0&-1&0&0&0&0\\0&0&-1&0&0&0
 \end{pmatrix}.
\end{equation}
Writing its rows as \(c_0,\ldots,c_5\), one has
\(c_jB_6=c_{j+1}\) for \(0\le j<5\) and
\(\det C=373248/125\).  These are identities in the full six-state
system.

For completeness, the angular elimination used below is finite and
unambiguous.  At a finite puncture \(p\), a leading pole row
\(v/(s-p)^m\), \(m\ge2\), is removed by subtracting the covariant
derivative of
\[
 \frac{v(R_p-(m-1)I)^{-1}}{(s-p)^{m-1}}.
\]
At infinity a polynomial leading row \(v s^d\) is removed by subtracting
the derivative of
\(v((d+1)I+R_0+R_1+R_\lambda)^{-1}s^{d+1}\).
The eigenvalues of \(R_0\) are \(1/6,1/2,5/6\), each repeated three
times; those of \(R_1,R_\lambda\) are \(0\) and \(-3/2\);
those of \(R_0+R_1+R_\lambda\) are
\(-1/6,-1/2,-5/6\), each repeated three times.  Every inverse just used
therefore exists.  A final constant primitive removes the simple zero
residue because \(R_0\) is invertible.  Finite primitives decay at
infinity and create only simple poles at other finite punctures;
polynomial primitives also create only simple finite poles.  Thus each
step strictly lowers the relevant higher pole or polynomial degree and
the procedure terminates.  The same leading-coefficient argument proves
uniqueness of the remaining form
\(v_1/(s-1)+v_\lambda/(s-\lambda)\).

Applying this stated procedure to \eqref{eq:appendix-Q} gives a rational
row \(\eta_S\) and six-row \(G(x)\) satisfying
\begin{equation}\label{eq:appendix-angular-source}
 \mathfrak q=\operatorname{NF}(G)+\partial_s\eta_S+\eta_S\mathcal M_s.
\end{equation}
The listing below constructs and retains the primitive;
its variable \texttt{eta}, divided by the coefficient-field
scaling \(\texttt{scale}=x^5(9x-25)(99x-50)^3\), is \(\eta_S\). It verifies this
entire nine-entry identity and the factorization of both residues in
\eqref{eq:appendix-NF}.  No unspecified angular witness is needed.

\subsection{The outer primitive and all its coefficients}

Use the six-coordinate solution column \(\mathcal V\) obtained by
multiplying \(\vv(\rho x^3)\) by \(U=\sqrt{1-x}\).
Thus \(\mathcal V=U\vv(\rho x^3)\), and its ordinary \(x\)-connection is
\[
 M_x=\frac{3B_6(\rho x^3)}x-\frac{I}{2(1-x)}.
\]
Since \(\dd U=-\dd x/(2U)\), the source row and three normalized period rows
in this frame are
\[
 r_S=-\frac{G}{2(1-x)},\qquad
 H_j=-\frac{c_j+6c_{j+1}}{2(1-x)},\qquad 0\le j\le2.
\]
Here \(D_\lambda\) denotes the coefficient action corresponding to
\(\theta_\lambda\); the full factor \(1+6D_\lambda\) is retained.  Define
\[
 (d_0,d_1,d_2)=\left(\frac1{375},\frac{14}{375},\frac{532}{3375}\right).
\]
An actual rational primitive is
\begin{equation}\label{eq:appendix-xi}
 \begin{aligned}
 \xi(x)&=\sum_{k=1}^{15}A_k\phi_k(x),\\
 (\phi_1,\ldots,\phi_6)&=
 \bigl(x^{-1},x^{-2},x^{-3},x^{-4},(99x-50)^{-1},(99x-50)^{-2}\bigr),\\
 (\phi_7,\ldots,\phi_{15})&=(1,x,\ldots,x^8).
 \end{aligned}
\end{equation}
All ninety rational entries of the rows \(A_k\) appear literally in
\texttt{COEFFICIENTS} in the listing.  No coefficient is a variable
to be solved for.  With that witness the exact identity is
\begin{equation}\label{eq:appendix-outer-identity}
 r_S-d_0H_0-d_1H_1-d_2H_2=\xi'+\xi M_x.
\end{equation}
It can be checked by clearing denominators in six rational functions.
Neither existence of a primitive nor a claim about a computed rank is
a premise of this equality.

The angular primitive enters the two-form identity with its wedge sign:
the exact one-form \(-\eta_S\,\dd U\), paired with the product solution
column, contributes
\((\partial_s\eta_S+\eta_S\mathcal M_s)\dd U\wedge\dd s\).
The outer term is the ordinary derivative of \(\xi\mathcal V\).
Hence \eqref{eq:appendix-angular-source} and
\eqref{eq:appendix-outer-identity} are the actual primitives required
for the comparison.  Their poles at \(99x-50=0\) remain present; the contour
truncation and endpoint argument in the main text justify applying the
identities to the selected literal contour.

\subsection{A complete exact verification recipe}

The following Python listing uses only exact SymPy arithmetic.  It is a
self-contained transcription of the finite identities: it reads no
previous certificate or saved output, executes no solver, and contains
the complete coefficient witness.  Equality of rational functions is
tested by exact polynomial cancellation.  Its source is obtained from
\eqref{eq:appendix-beta}, not inserted as a fitted six-row.  During
angular reduction, multiplying by
\(x^5(9x-25)(99x-50)^3\) is only a coefficient-field scaling and is
undone in the result.  For generic \(x\) and \(\lambda\ne0,1\), all
steps are defined; the resulting rational identities consequently hold
throughout their domains by clearing denominators.

% The complete listing is appended by the owned document generator.
% Complete exact certificate listing.
\begin{lstlisting}[language=Python,basicstyle=\ttfamily\scriptsize,breaklines=true,columns=fullflexible,keepspaces=true,showstringspaces=false]
"""Standalone exact transcription check for the printed rational appendix.

This file supplies its own literal inputs and never reads or executes an old
certificate, solver, or output.  No period, target fit, or rank is an input.
The printed 90 coefficients constitute a witness, not unknowns to be solved.
"""
import sympy as S
from sympy.polys.fields import field
import time

started = time.perf_counter()
X, Z, LAM = S.symbols('x s lambda')
rho = S.Rational(729, 15625)
I = S.eye(3)
J = S.Matrix([[0, 1, 0], [0, 0, 1], [0, 0, 0]])
w = S.Matrix([[S.Rational(5,72), S.Rational(23,36), S.Rational(3,2)]])
E = S.zeros(3)
E[2,:] = w
R0 = S.kronecker_product(J,I) - S.kronecker_product(I,J-E)
R1 = -S.kronecker_product(E,I)
Rl = -S.kronecker_product(I,E)
Rt = R0 + R1 + Rl
B6 = (J + LAM*E/(1-LAM)).row_join(-LAM*E/(1-LAM)).col_join(
    (-E/(1-LAM)).row_join(J + LAM*E/(1-LAM)))
C = S.Matrix([
    [S.Rational(49456,125),S.Rational(153648,125),S.Rational(113472,125),
     S.Rational(72,5),-S.Rational(3312,25),S.Rational(113472,125)],
    [-S.Rational(1576,25),-S.Rational(4608,25),-S.Rational(3312,25),
     0,S.Rational(72,5),-S.Rational(3312,25)],
    [S.Rational(46,5),S.Rational(108,5),S.Rational(72,5),0,0,S.Rational(72,5)],
    [-1,0,0,0,0,0],[0,-1,0,0,0,0],[0,0,-1,0,0,0]])

def cancel_matrix(A):
    return A.applyfunc(S.cancel)

def nf(v):
    return (S.kronecker_product(w,v[:,:3])/(Z-1)
            + S.kronecker_product(v[:,3:],w)/(Z-LAM))

N = lambda q: J + q*E/(1-q)
gauge = S.Matrix([[1,0,0],[0,LAM,0],
                 [5*LAM/(72*(1-LAM)),LAM**2/(2*(1-LAM)),2*LAM**2]])
pq,cq = (1-S.Rational(3,2)*LAM)/(LAM*(1-LAM)),5/(144*LAM*(1-LAM))
ao = S.Matrix([[0,1,0],[2*cq,-pq,2],[0,cq,-2*pq]])
jc = S.Matrix([[-5*LAM/36,-LAM/2,1-LAM],[-LAM/2,LAM-1,0],[1-LAM,0,0]])
tq = S.Matrix([[-5*LAM/72,0,2*LAM**2*(1-LAM)],
               [-LAM/2,LAM*(LAM-1),0],[1-LAM,0,0]])
assert cancel_matrix(gauge.diff(LAM)+gauge*ao-N(LAM)*gauge/LAM) == S.zeros(3)
assert cancel_matrix(jc*gauge-tq) == S.zeros(3)
Ms = S.kronecker_product(N(Z),I)/Z - S.kronecker_product(I,N(LAM/Z))/Z
Ml = S.kronecker_product(I,N(LAM/Z))
assert cancel_matrix(Ms-R0/Z-R1/(Z-1)-Rl/(Z-LAM)) == S.zeros(9)
for j in range(6):
    v = S.eye(6)[j,:]
    p = S.zeros(1,9) if j < 3 else -Z*S.kronecker_product(v[:,3:],w)/(Z-LAM)
    assert cancel_matrix(LAM*nf(v).diff(LAM)+nf(v)*Ml-nf(v*B6)
                         -p.diff(Z)-p*Ms) == S.zeros(1,9)
e00 = S.eye(9)[0,:]
anchor = e00*R0.inv()
assert cancel_matrix(e00/Z-nf(C[0,:])-anchor*Ms) == S.zeros(1,9)
for j in range(5):
    assert cancel_matrix(C[j,:]*B6-C[j+1,:]) == S.zeros(1,6)
assert C.det() == S.Rational(373248,125)

# The rational beta source, with beta_N = T_U dU + U*T_s ds.
x, s, lam = X, Z, rho*X**3
m = 27*x*(80-99*x+27*x**2)/1000
a = 27*x*(9*x-16)/1000
M = 25-9*x
D = x**4*M*(99*x-50)**2
f = s**2-2*m*s+lam
a2 = 81*x**2-198*x+80
a3 = 243*x**3-1053*x**2+1266*x-400
a6 = (59049*x**6-511758*x**5+1724085*x**4-2860596*x**3
      +2441988*x**2-1011200*x+160000)
k,jj,vv,cc = map(S.Rational, ['1953125/209952','244140625/944784',
                             '3906250/6561','30517578125/8503056'])
LP = S.Matrix([[0,0,-k*(9*x-25)*a2*a3],
               [0,0,jj*(9*x-25)*a2],
               [k*(9*x-25)*a2*a3,-jj*(9*x-25)*a2,0]])
RP = S.Matrix([[0,0,k*a6],[0,vv*x*(99*x-50),-jj*a3],
               [k*a6,-jj*a3,cc]])

# These temporary matrices are the displayed beta numerators.
# Their common scalar denominator is D*f.
Ts = cancel_matrix(2*LP*(s-m)+2*a*M*RP)
TU = cancel_matrix((-2*(1-x)*LP*S.diff(f,x)-2*RP*(
     -2*a*(1-x)*M*S.diff(m,x)+(s-m)*(
     (a-2*(1-x)*S.diff(a,x))*M+9*(1-x)*a))))

# The beta numerators are polynomials in s over Q(x).
# Clear the known connection denominator, divide by f^2 exactly, and only
# then create rational rows.  This preserves all nine source entries.
F, x = field('x',S.QQ)
K, s = field('s',F.to_domain())
z = K.zero
lam = rho*x**3
lift = K.ground_new
parse = lambda q: K.from_expr(q)
constant = lambda A: [[parse(A[i,j]) for j in range(A.cols)] for i in range(A.rows)]
ts_fraction,tu_fraction = constant(Ts),constant(TU)
assert all(q.denom.degree() == 0 for row in ts_fraction+tu_fraction for q in row)
ts = [[q.numer.mul_ground(F.one/q.denom.LC) for q in row] for row in ts_fraction]
tu = [[q.numer.mul_ground(F.one/q.denom.LC) for q in row] for row in tu_fraction]
P = K.ring
sp = s.numer
pp = P.one
zz = P.zero
u = 1-x
mfield = F.from_expr(m)
dlogD = F.from_expr(S.diff(D,X)/D)
fpoly = sp**2-2*mfield*sp+lam

def dx_polynomial(p):
    return P.from_dict({e:c.diff(x) for e,c in p.items()})

fx,fs = dx_polynomial(fpoly),fpoly.diff(sp)
angular_den = sp**2*(1-sp)*(sp-lam)
asbar = [[zz,sp*(1-sp),zz],
         [S.Rational(5,72)*pp,-pp+S.Rational(3,2)*sp,2*sp*(1-sp)],
         [zz,S.Rational(5,144)*pp,-2*pp+3*sp]]
cbbar = [[zz,lam*(sp-lam),zz],
         [S.Rational(5,72)*sp**2,-sp**2+S.Rational(3,2)*lam*sp,2*lam*(sp-lam)],
         [zz,S.Rational(5,144)*sp**2,-2*sp**2+3*lam*sp]]
scale = lift(x**5*(9*x-25)*(99*x-50)**3)
scale_over_D = -x*(99*x-50)
alpha = [[z]*3 for _ in range(3)]
for i in range(3):
    for j in range(3):
        scalar_part = (-ts[i][j]*fpoly
            +2*u*(dx_polynomial(ts[i][j])*fpoly-ts[i][j]*fx
                  -dlogD*ts[i][j]*fpoly)
            +tu[i][j].diff(sp)*fpoly-tu[i][j]*fs)
        numerator = (angular_den*scalar_part
            -sp*(sp-lam)*sum((asbar[i][k]*tu[k][j] for k in range(3)),zz)*fpoly
            +(1-sp)*sum((tu[i][k]*cbbar[j][k] for k in range(3)),zz)*fpoly
            -(6*u/x)*sp*(1-sp)*sum((ts[i][k]*cbbar[j][k] for k in range(3)),zz)*fpoly)
        quotient,remainder = numerator.div(fpoly**2)
        assert remainder == zz
        alpha[i][j] = K.new(quotient.mul_ground(scale_over_D),angular_den)
print('All nine source entries pass exact zero-remainder f^2 division.', flush=True)

def metric_gauge(q):
    return [[-5*q/72,z,2*q*q*(1-q)],
            [-q/2,q*(q-1),z],[1-q,z,z]]

b = lift(lam)/s
left,right = metric_gauge(s),metric_gauge(b)
Q = [[sum((left[i][a]*alpha[a][c]*right[j][c]
           for a in range(3) for c in range(3)),z)
      for j in range(3)] for i in range(3)]
print('Source transcribed from the complete beta formulas.', flush=True)

r0,r1,rl = map(constant,[R0,R1,Rl])
ms = [[r0[i][j]/s+r1[i][j]/(s-1)+rl[i][j]/(s-lift(lam))
       for j in range(9)] for i in range(9)]

def row_times(v,A):
    return [sum((v[i]*A[i][j] for i in range(len(v))),z)
            for j in range(len(A[0]))]

def covariant(v):
    return [v[j].diff(s)+q for j,q in enumerate(row_times(v,ms))]

def evaluate(q,a):
    def ev(p):
        return sum((c*(a**e[0] if e[0] else F.one) for e,c in p.items()),F.zero)
    return lift(ev(q.numer)/ev(q.denom))

def pole_order(q,a):
    p,linear,n = q.denom,(s-lift(a)).numer,0
    while q:
        quotient,remainder = p.div(linear)
        if remainder:
            break
        p,n = quotient,n+1
    return n

q = [Q[i][j] for i in range(3) for j in range(3)]
r, eta = q[:],[z]*9

def subtract(p):
    global r,eta
    exact = covariant(p)
    r = [r[j]-exact[j] for j in range(9)]
    eta = [eta[j]+p[j] for j in range(9)]

for a,residue in [(F.one,R1),(lam,Rl),(F.zero,R0)]:
    while True:
        order = max(pole_order(t,a) for t in r)
        if order < 2:
            break
        lead = [evaluate(t*(s-lift(a))**order,a) for t in r]
        inv = constant((residue-(order-1)*S.eye(9)).inv())
        subtract([t/(s-lift(a))**(order-1) for t in row_times(lead,inv)])
while True:
    polynomials = [t.numer.div(t.denom)[0] for t in r]
    degree = max(p.degree(0) if p else -1 for p in polynomials)
    if degree < 0:
        break
    lead = [lift(p.get((degree,),F.zero)) for p in polynomials]
    inv = constant(((degree+1)*S.eye(9)+Rt).inv())
    subtract([t*s**(degree+1) for t in row_times(lead,inv)])
subtract(row_times([evaluate(t*s,F.zero) for t in r],constant(R0.inv())))
u = [evaluate(t*(s-1),F.one) for t in r]
v = [evaluate(t*(s-lift(lam)),lam) for t in r]
normal = [u[j]/(s-1)+v[j]/(s-lift(lam)) for j in range(9)]
assert all(not (r[j]-normal[j]) for j in range(9))
exact = covariant(eta)
assert all(not (q[j]-normal[j]-exact[j]) for j in range(9))
g = [u[6+j]/parse(w[0,2])/scale for j in range(3)]
g += [v[3*j+2]/parse(w[0,2])/scale for j in range(3)]
assert all(not (u[3*i+j]-parse(w[0,i])*g[j]*scale)
           and not (v[3*i+j]-g[3+i]*parse(w[0,j])*scale)
           for i in range(3) for j in range(3))
print('Angular source identity established with its rational primitive.', flush=True)

# Literal 15 by 6 coefficient table: basis order is specified below.
COEFFICIENTS = '''
105800/459459 867560/459459 103229200/459459 -169280/51051 -1388096/51051 -2268352/51051
0 0 -472485400/459459 0 0 0
0 0 338560000/196911 0 0 0
0 0 -4232000000/4135131 0 0 0
-78890805485931260/886402053948411 -1227125311357021190/886402053948411 -70427029255928891500/6204814377638877 -607873567055897705440/74327471429736107583 -298973396640778498112/8258607936637345287 -145365088004028281168/635277533587488099
-231954922267076000/229807939912551 -331597299956488000/76602646637517 -108998157631733600/5892511279809 -39913346954892419000/917623104070816143 28252939480465998800/101958122674535127 -14427468858797595200/5997536627913831
-44128649359258592/31024071888194385 623049674458568524/31024071888194385 6660655461825680912/31024071888194385 57598947351747925330/10618210204248015369 370356988491186478316/8258607936637345287 422172179336846381488/5899005669026675205
7410960991604618219/7834361587927875000 -2244024947258679397/230422399644937500 -8368444588355333518/75330399883921875 -1674067760127596905471/469239087308940073125 -7571431721389433993041/260688381838300040625 -60217890176084166521579/1303441909191500203125
-1014783336398633/79134965534625000 99347762374288343/39567482767312500 226386615454772638/9891870691828125 1399948922915894608/4739788760696364375 6901571092120745041/2633215978164646875 59461141531221275693/13166079890823234375
-1051211136/19184234069 -440067361152/479605851725 -18358156683264/2398029258625 1031881871061389312/1196916353711203125 258130738760714016/40711440602421875 20750847913747165488/1994860589518671875
1029377673/34880425580 430927242561/872010639500 4494215281788/1090013299375 -250759381828805733/558002962103125000 -14957478836244716238/4533774067087890625 -24572653010256580956/4533774067087890625
-16224867/3170947780 -6792198219/79273694500 -70837018452/99092118125 300193629643629927/3297290230609375000 1380547922050429344/2060806394130859375 2261266906791518928/2060806394130859375
0 0 0 -223303109482368/7493841433203125 -8431993935111168/37469207166015625 -13467359214383616/37469207166015625
0 0 0 218665144746549/13625166242187500 2064215740879206/17031457802734375 3296911156760547/17031457802734375
0 0 0 -3446560950471/1238651476562500 -32535799768674/1548314345703125 -51965324712513/1548314345703125
'''
coefficients = [[F.from_expr(S.Rational(t)) for t in line.split()]
                for line in COEFFICIENTS.strip().splitlines()]
assert len(coefficients) == 15 and all(len(row) == 6 for row in coefficients)
basis = [x**(-p) for p in range(1,5)]
basis += [(99*x-50)**(-p) for p in range(1,3)] + [x**p for p in range(9)]
xi = [sum((coefficients[k][j]*basis[k] for k in range(15)),F.zero)
      for j in range(6)]
mx = [[3*F.from_expr(B6[i,j].subs(LAM,lam.as_expr()))/x
       -(F.one/(2*(1-x)) if i == j else F.zero)
       for j in range(6)] for i in range(6)]
d = list(map(S.Rational,['1/375','14/375','532/3375']))
h = [[-F.from_expr(C[k,j]+6*C[k+1,j])/(2*(1-x))
      for j in range(6)] for k in range(3)]
for j in range(6):
    residual = (-g[j].as_expr()/(2*(1-X))
                -sum(d[k]*h[k][j].as_expr() for k in range(3))
                -xi[j].diff(x).as_expr()
                -sum((xi[i]*mx[i][j]).as_expr() for i in range(6)))
    assert S.cancel(residual) == 0
assert [-81*t for t in d] == list(map(S.Rational,['-27/125','-378/125','-1596/125']))
print('All finite printed identities hold over the rational function fields.')
print('Elapsed seconds:', round(time.perf_counter()-started,3))
\end{lstlisting}


\section{Provenance and review boundary}\label{app:provenance}
The finite rational inputs were transcribed from the frozen Row12 proof
package associated with author commit
\texttt{a17703b4e69a5d879d71ebd37fcaba32af45be0c}
and strict-audit package
\texttt{7b71f2c0f1c72798dda3e8759da4465b489366c0}.
In that package the relevant inputs are the source-lift proof and
\texttt{s03/LIFT\_CHECK.output.json}, the six-coordinate proof
\texttt{s08/PROOF.md}, and the complete outer coefficients
\texttt{s14/OUTER\_ROW\_CHECK.output.json}.
The finite formulas and verification procedure used in this document are
printed in Appendix~\ref{app:certificates}, so those historical files are
provenance references rather than additional compilation inputs.

The analytic scalar route was organized in
\texttt{SCALAR\_INTEGRAL\_ROUTE.md} and checked together with the literal
coefficient and moment bridge in the additional scalar review performed
7 October 2026. The signed Gauss calculation uses the corrected actual
derivatives recorded in the completed signed-Gauss Lean support package;
the earlier unchecked draft is not a premise of this proof.
The document-level review, its checked source version and any corrections
are recorded next to this source in \texttt{DOCUMENT\_REVIEW.md}.

The evidential status is [ARG] for the analytic derivation and [ONE] for
the exact document-transcription checks. Human acceptance is Pending.
The connected Lean formalization was completed on 7 October 2026. Its
unconditional declarations are \texttt{Row12.row12\_hasSum},
\texttt{Row12.row12\_summable} and \texttt{Row12.row12\_tsum}.
The final statement and axiom outputs, all five verification checks and the
fresh full-result review are retained in the separate
\path{C5/row12_lean_2026-10-06} package. That formal result covers exactly
the scalar equality proved here.
\end{document}
